\chapter{Conclusions}
\label{sec:conclusion}

Small satellite technology has increased the availability of remote sensing data such as hyperspectral imagery. The image sensor on-board the \acrshort{hypso} missions developed at \acrshort{NTNU} capture images of 120 spectral bands, producing large amounts of data that must be downlinked with limited bandwidth. On-board image compression is needed to effectively handle image data, and is typically implemented as hardware accelerators running on flexible \acrshort{FPGA}s.

The CCSDS 123.0-B-2 compression algorithm (issue 2) supports near-lossless compression capable of high compression performance. The complexity of the algorithm, and intersample data dependencies, presents challenges for hardware accelerators that want to optimize hardware resource utilization, throughput and power. Three main approaches have been attempted in the literature to achieve high throughput implementations. \citeauthor{basconesRealTimeFPGAImplementation2022} suggest using a novel sample processing order to interleave data dependant calculations by up to $N_z$ samples. Using VHDL \acrshort{RTL} the approach reaches a throughput of 285 MS/s, but at a large area footprint of 16458 \acrshort{LUT}s and 15707 \acrshort{FF}s. \citeauthor{chatziantoniouScalableDataRateNearLossless2022a} and \citeauthor{valsesiaHighThroughputOnboardHyperspectral2019} suggest using prequantization, in place of the in-loop quantizer of the standard. This approach matches the throughput of \citeauthor{basconesRealTimeFPGAImplementation2022}, but at a penalty in compression performance and even greater area footprint. \citeauthor{sanchezReducingDataDependencies2022} suggest a mathematically equivalent method for calculating the sign of the double-resolution prediction error needed for updating the adaptive weight vector of the prediction loop. This, combined with using reduced prediction mode, narrow local sums and the \acrshort{BIL} processing order, allows implementation of a pipeline that is limited by the weight update calculation. Using this approach, \citeauthor{vorhaugHyperspectralImageCompression2024} designs a compression core reaching a throughput of 110 MS/s with the smallest area footprint at 5577 \acrshort{LUT}s and 2635 \acrshort{FF}s. The design is successfully tested on-board the \acrshort{hypso}-1 satellite. To alleviate the critical path of the mathematical optimization approach, \citeauthor{jiaRemovalFeedbackLoop2025} suggest a novel weight update scheme, dubbed delayed weight updates, reaching a throughput of 348 MS/s for the predictor module.

This work integrates the delayed weight update scheme into the accelerator designed by \citeauthor{vorhaugHyperspectralImageCompression2024}. The new pipeline reaches 202 MS/s for the Zynq ZU7EV \acrshort{FPGA} with a minimal increase in area for a total of 6349 \acrshort{LUT}s and 3480 \acrshort{FF}s. In isolation, the predictor module reaches 351 MS/s, making it the fastest and smallest implementation in the open literature. The design is successfully tested on-board the \acrshort{hypso}-2 satellite, with a throughput of 92 MS/s on the Zynq-7030. Future improvements to the design should focus on increasing the throughput of the limiting hybrid encoder module.

The new weight update scheme deviates from the issue 2 standard, requiring custom software for decompression. This works presents an open-source, issue 2 compliant decompressor based on the verification model by \citeauthor{vorhaugHyperspectralImageCompression2024} with optional support for delayed weight updates. The new predictor module written in C++ achieves a factor seven improvement in execution time. Further improvements in execution time can be gained in the encoder module and by optimizing memory accesses. The verification module is used to verify the hardware accelerator and decompress its outputs. For \acrshort{hypso}-2 images, the penalty in compression performance due to delayed weight updates is less than one percent, an acceptable tradeoff for the increase in throughput.

